\section{Introduction}\label{sec:introduction}
Let $(X,B)$ be a projective log canonical pair with rational boundary.
Its log canonical divisor $K_X+B$ records the canonical class together
with the boundary. A rational Cartier divisor is \emph{nef} if its degree
on every curve is nonnegative. It is \emph{semiample} if a positive Cartier
multiple is generated by global sections. Such a multiple defines a morphism,
so semiampleness turns numerical positivity into a fibration described by
sections. The log abundance conjecture asks whether a nef log canonical
divisor is semiample. We give a positive answer in characteristic zero.

\begin{theorem}[Log abundance]\label{thm:main}
Let $(X,B)$ be a normal projective log canonical pair over an algebraically
closed field of characteristic zero. Assume that $B$ is an effective rational
divisor and $K_X+B$ is $\Q$-Cartier. If $K_X+B$ is nef, then some positive
Cartier multiple of it is generated by global sections.
\end{theorem}

The proof proceeds by dimension induction over $\C$. Its inductive assertion
is stronger than the displayed theorem: every projective lc pair with effective
real boundary and pseudo-effective real Cartier adjoint has a good log minimal
model, on which the adjoint is semiample. Smooth canonical nonvanishing asks whether a pseudo-effective canonical
class on a smooth variety has a nonzero pluricanonical section. We
establish this statement in each dimension before
that good-model assertion is used in the same dimension. At the end, descent
of the evaluation map of an actual Cartier line bundle gives the theorem over
every algebraically closed field of characteristic zero.

We write $\kappa(L)$ for the Iitaka dimension of a rational divisor $L$,
with value $-\infty$ if all positive integral systems are empty.
The consequences fall into three groups. Good models, rational-linear
triviality when $\kappa=0$, and gluing and extension across boundaries are
treated in Section~\ref{sec:consequences}. Finite generation of the algebras
of rounded pluricanonical sections is proved in
Section~\ref{sec:canonical-rings}, both absolutely in characteristic zero
and for proper morphisms over algebraic bases over $\C$.
Section~\ref{sec:classification} relates canonical nonvanishing to rational
curves, cotangent tensors, fundamental groups, and quasi-projective covers.
These sections give the full statements, additional hypotheses, and
derivations from the work of Fujino--Gongyo, Boucksom--Demailly--P{\u a}un--Peternell,
Lazi{\'c}--Peternell, Brunebarbe--Campana, Campana, Taji, and
Claudon--H\"oring--Koll\'ar.

There is also a relative analytic application of Fujino's theorem:
Section~\ref{sec:analytic-abundance} treats projective surjective morphisms
of normal complex analytic varieties with rational lc adjoint nef over
the whole base. It obtains a relatively generated Cartier multiple near
each compact base subset; the multiple may depend on that subset.
This is distinct from the compact K\"ahler symmetric-pluriform criterion
in Section~\ref{sec:classification}; the cotangent-invariant and
fundamental-group consequences there remain projective.

\subsection*{Historical development}
The threefold proofs show why nonvanishing and boundary geometry are
central to abundance. Miyaoka proved the numerical-dimension-one case
by studying an effective pluricanonical divisor and its infinitesimal
neighborhoods after suitable coverings \cite{Miyaoka1988}. Kawamata proved
the remaining numerical-dimension-two case for minimal threefolds
\cite[Theorem~3.1]{Kawamata1992}. His Section~4 gives an alternative proof of
Miyaoka's case through finite covers, Hodge theory and compatible
infinitesimal deformations. Keel, Matsuki and
McKernan established the logarithmic theorem
\cite{KeelMatsukiMcKernan1994,Matsuki2003Correction,
KeelMatsukiMcKernan2004Correction}.
Fujino extended abundance to semi-log-canonical threefolds
\cite{Fujino2000}, where sections on the normalization must agree along
the conductor before they define sections on the original space. This
compatibility is also needed in higher-dimensional induction: the
coefficient-one boundary can be reducible even when the ambient variety
is normal.

In higher dimensions, Birkar, Cascini, Hacon and McKernan established
minimal models for klt pairs in the big-boundary or log-general-type range
\cite{BCHM2010}. Their finite-generation theorem for
projective rational klt pairs does not require bigness
\cite[Corollary~1.1.2]{BCHM2010}. The rational lc finite-generation
consequence below goes beyond that klt setting.
Two distinctions remain essential on the boundary
of the positive cone. First, nonvanishing produces sections but does not
by itself assert that a nef adjoint is semiample. Hashizume reduces
real-boundary log canonical nonvanishing and ordinary minimal-model
existence through dimension $n$ to smooth canonical nonvanishing in
dimension $n$ \cite[Theorem~1.4]{Hashizume2018}. Our induction must then
make these models good. For positive Kodaira dimension, the required
step is supplied by lower-dimensional good models
\cite[Lemma~3.5]{FujinoGongyoRings2014}. Second, Kodaira dimension zero
and numerical dimension zero are different hypotheses. Gongyo proves
that a numerically trivial rational log canonical adjoint is
$\Q$-linearly trivial \cite[Theorem~1.2]{Gongyo2013}. In the
zero-Iitaka argument below, numerical triviality forces an already
available effective representative to vanish.

\begingroup\clubpenalty=10000
Recent work of Liu and Xu treats numerical dimension at most one.
Their Theorem~5.1 gives good minimal models for projective log canonical
pairs of dimension at most five with nonnegative invariant Kodaira
dimension and numerical dimension at most one; their Theorem~5.4 gives
a reduction in that numerical range to smooth nonvanishing
\cite{LiuXu2025}. The induction here has no numerical-dimension restriction.
The ordinary-pair specialization of the companion generalized-minimal-model
paper \cite[Corollary~1.2]{OpenAIMinimalModels2026} concerns existence of a
nef model. It is not an input to this good-model induction. After abundance
has been proved, we use that ordinary-pair result in the absolute
finite-generation argument of Section~\ref{sec:canonical-rings}.
\par\endgroup

\subsection*{The two tasks in the induction}
Assuming good models in dimensions below $n$, we first prove smooth
canonical nonvanishing in dimension $n$, and then make the resulting
log minimal models good. Both tasks use a criterion for a \emph{signed}
representative on a reduced boundary. The lower-dimensional hypothesis
supplies its whole-boundary semiampleness assumption.
We describe this common ingredient first.
For a nef divisor $L$, write $\nu(L)=\max\{j:L^jH^{n-j}>0\}$,
where $H$ is ample and $n=\dim X$.
The criterion is
\[
 \begin{gathered}
 L=K_X+D\text{ nef},\qquad L-cD\text{ pseudo-effective }(c>0),\\
 L\sim_{\Q}\sum_i a_iD_i,\qquad L|_D\text{ semiample}
 \quad\Longrightarrow\quad\kappa(L)=\nu(L).
 \end{gathered}
\]
Here $(X,D)$ is projective $\Q$-factorial dlt, and
$D=\sum_iD_i$ is the decomposition of the reduced boundary into its prime
components. The rational coefficients $a_i$ may be positive, negative, or
zero. Semiampleness is required on the entire
reduced scheme $D$. Normalization and conductor gluing therefore belong to
the induction, as in the slc interface of Fujino--Gongyo
\cite[Theorem~1.5]{FujinoGongyo2014}.
For the good-model task, positive Kodaira dimension is handled by
lower-dimensional good models; the signed criterion addresses the
remaining zero-Iitaka-dimension case in
Proposition~\ref{prop:inductive-reduction}.

The signed criterion is proved geometrically. A generated multiple of
$L|_D$ defines a morphism from $D$ to projective space. Over a general point
of the image of a positive-coefficient component, a root construction
separates that component from the negative and coefficient-zero boundary.
A filtered Hodge-module calculation lifts local parameters on this image,
together with a parameter transverse to the lifted boundary, through all
infinitesimal neighborhoods. Fixing the lifted base parameters and varying
the transverse one deforms a boundary fiber into a projective subvariety
disjoint from $D$. Varying the fiber gives a dominating family on which $L$
is trivial. Section~\ref{sg:descent} explains the final geometric step:
these subvarieties bound the dimension of the nef reduction, and vertical
descent identifies $L$ with the pullback of a big divisor on its base.

This formal
deformation method has a direct precedent in Miyaoka's construction;
Miyaoka credits Reid with the suggestion to analyze a pluricanonical
divisor \cite[p.~220]{Miyaoka1988}. Boundary extension provides a second
related approach. Demailly, Hacon and P\u{a}un use a singular-metric
refinement of Ohsawa--Takegoshi extension to extend pluricanonical sections
in the purely log terminal setting
\cite[Theorem~1.7 and Corollary~1.8]{DemaillyHaconPaun2013}.
Their nef corollary requires an effective representative whose support
contains the coefficient-one divisor and lies in the boundary. Chan and Choi subsequently removed the containment of that representative
in the boundary \cite[Theorem~1.6.1]{ChanChoi2021}.
Here the representative may have both signs, and compatibility is required
across the whole reducible boundary.
We therefore construct compatible formal lifts on the whole reduced
boundary. Semiampleness on that boundary is supplied by the normalization
theorem of Fujino and Gongyo
\cite[Theorem~1.5]{FujinoGongyo2014}.

The construction combines familiar tools with an additional compatibility
argument. Normalized cyclic covers and their eigenspaces follow the
formalism of Esnault and Viehweg \cite[Section~3]{EsnaultViehweg1992}.
Deligne's logarithmic Hodge theory supplies the underlying framework
\cite{Deligne1971}. On the singular boundary, we use Saito's projective
strictness and graded de Rham vanishing, together with his comparison
with the Du Bois complex
\cite[Theorem~2.14 and Proposition~2.33]{Saito1990}
\cite[Theorem~0.2 and Corollary~0.3]{Saito2000}.
The proof must still identify the obstruction to lifting and show that
these ordinary-variety results apply to the covering construction.

\paragraph{Smooth nonvanishing.}
To prove smooth nonvanishing, suppose instead that a counterexample exists.
The zero-boundary case of logarithmic Iitaka subadditivity, stated below
as Theorem~\ref{asm:iitaka}, first excludes positive irregularity through
the Albanese map. This is its only use in the proof.
The geometric preparations rule out two ways it could be covered by
lower-dimensional subvarieties: families that carry a positive current
pulled back from a smaller space, and covering families of curves with bounded genus
and bounded degree after normalizing the nef divisors used in the
construction to a fixed small positive volume. The first exclusion uses positive currents and algebraic webs.
The second uses an ascending chain condition for the singularity
multiplicities of one fixed current, measured by Lelong numbers at valuations
of bounded discrepancy.
The signed criterion has a second role here: it forces the full reduced
support of any signed canonical representative to have big logarithmic
adjoint, as proved in Proposition~\ref{prop:signed-alternative}. This will
control the multisections in the two-slot construction below.

These exclusions produce two incompatible estimates for vanishing at a
moving point. At a very general point, the first estimate bounds the
vanishing order of every nonzero section of a Cartier multiple $kP$ of a suitable
nef divisor $P$ by $k$ times a small normalized constant.
To obtain sections with stronger jet separation, we use the
projective bundle associated to the sum of the two pullbacks of a line bundle
to the product of a smooth model with itself. Its two summands allow the
argument to move in either factor.

Failure of the required jet separation would produce a moving base
component. Slice estimates and the curve exclusions leave only components
that project generically finitely onto both base factors, hence give
correspondences between them. Their branch divisors cannot sweep either
factor; a fixed branch complement and bounded degree therefore leave only
finitely many covers. After fixing the covers, the correspondences give
a family of birational maps whose graphs dominate their product. Hanamura's
birational-group theorem then makes a fixed cover birational to an abelian
variety \cite[Theorems~2.1--2.2]{Hanamura1988}. The ramification and
positive-current argument in Section~\ref{sec:exclusions} would give a
pluricanonical section on the counterexample, again a contradiction.

Once the projective bundle, its divisors, and the
required finite collection of jet sections are fixed, reduction modulo a
large prime gives an evaluation map on the product of two copies of the
Frobenius-twisted model. The map has full generic rank. Its restriction
to the ordinary diagonal is controlled by sections on the Frobenius
thickening of the diagonal, where the two relative Frobenius maps agree.
A filtration of those sections gives a much smaller rank bound there.
The determinant is a section of an external product of two line bundles.
The first estimate, expressed by a fixed movable curve on a blowup,
bounds vanishing in each factor and hence the order of the determinant
at the selected diagonal point. The rank deficit gives a larger lower
bound on that order, a contradiction.

Theorem~\ref{common:finite-data-frobenius} states this incompatibility for
a fixed smooth variety, fixed jet sections, and a fixed movable-curve bound.
Its proof opens Section~\ref{fr:section} and uses neither abundance nor
subadditivity. The geometric application follows in the same section,
using the fixed inputs supplied by the preceding curve exclusions and
jet constructions. The numerical subadjunction and moving-kernel results
needed for those constructions are proved in Section~\ref{sec:common-tools}.

The numerical part combines two established methods.  The argument that
follows a moving base component uses differentiation in the parameter
space, as in Ein--K\"uchle--Lazarsfeld
\cite[Proposition~2.3]{EinKuchleLazarsfeld1995}; we give the tracking and
adjunction estimates needed here.  After reduction to positive
characteristic, we compare ordinary jets with Frobenius jets using the
local containment recorded by Musta\c{t}\u{a}--Schwede
\cite[Equation~(2.8)]{MustataSchwede2014}.  The resulting evaluation map
is studied through Sun's canonical Frobenius filtration
\cite[Theorem~3.7]{Sun2008}, whose curve case is treated by
Joshi--Ramanan--Xia--Yu \cite[Section~5.3]{JoshiRamananXiaYu2006}.
We use its description by powers of the diagonal ideal due to
Kitadai--Sumihiro \cite[Definition~3.1 and Remark~3.2]{KitadaiSumihiro2008}.
Langer's instability estimate \cite[Corollary~2.5]{Langer2004} then
controls the slopes of the graded bundles.

\subsection*{The subadditivity input}
The all-dimensional argument uses the following completed companion result.
Its application occurs only in Lemma~\ref{ex:irregularity}.

\begin{theorem}[Logarithmic Iitaka subadditivity,
{\cite[Corollary~6.2]{OpenAIIitaka2026}}]
\label{asm:iitaka}
Let \(f:X\to Y\) be a surjective morphism with connected fibers between
smooth connected projective complex varieties. Let \(D_X,D_Y\) be
reduced effective simple normal crossing divisors, allowing zero
boundaries, such that
\[
  \Supp(f^*D_Y)\subseteq\Supp(D_X).
\]
For a very general smooth fiber \(F\), put \(D_F=D_X|_F\). Then
\[
 \kappa(X,K_X+D_X)\ge
 \kappa(F,K_F+D_F)+\kappa(Y,K_Y+D_Y).
\]
Here \(D_F\) is reduced with simple normal crossings and
\(K_F+D_F\sim(K_X+D_X)|_F\). The convention
\((-\infty)+b=-\infty\) applies also when \(b=-\infty\), and the zero
divisor on a point has Iitaka dimension zero.
\end{theorem}

The support inclusion allows additional components in \(D_X\), and
the morphism need not be smooth away from the boundary. The theorem
has no nefness, abundance, bigness, or good-model hypothesis. Its
zero-boundary case is used in the Albanese reduction of
Section~\ref{sec:exclusions}; this is the sole subadditivity input in
our proof. The separate characteristic-zero specialization
\cite[Corollary~6.3]{OpenAIIitaka2026} uses a geometric generic fiber.
We keep that field-and-fiber statement distinct from the complex,
very-general-fiber statement above. No fractional-boundary,
orbifold, nonprojective, Hodge-conjectural, or arithmetic extension
is assumed here.

The particular zero-boundary application in Lemma~\ref{ex:irregularity}
also follows from the external maximal-Albanese-dimension theorem of
Hacon, Popa and Schnell \cite[Theorem~1.1]{HaconPopaSchnell2017}:
the smooth base there maps generically finitely to a subvariety of an
abelian variety. Cao and P\u{a}un prove the underlying abelian-base
case \cite[Theorem~1.1]{CaoPaunAbelian2016}. These results cover that
application, not the full logarithmic statement above.

\paragraph{Reading order.}
Section~\ref{sec:induction} sets out the dimension induction.
Sections~\ref{sg:section}--\ref{sg:descent} prove the signed-boundary criterion:
first the root and adjunction construction, then infinitesimal lifting, and
finally compact fibers and descent. Section~\ref{sec:exclusions} derives the
geometric exclusions for a hypothetical smooth counterexample.
After the local tools of Section~\ref{sec:common-tools},
Section~\ref{sec:jets} constructs its scalar and two-slot jets.
Section~\ref{fr:section} proves the independent Frobenius comparison and
then constructs its fixed witnesses to establish smooth nonvanishing.
Sections~\ref{sec:completion} and~\ref{sec:consequences} finish the induction,
descend actual global generation, and record good models, linear triviality,
semi-log-canonical abundance, and supported dlt extension.
Section~\ref{sec:canonical-rings} then proves finite generation of rational
log canonical rings in the absolute and complex proper relative settings.
Section~\ref{sec:classification} combines canonical nonvanishing and
abundance with classical results on rational curves, cotangent tensors,
fundamental groups, and quasi-projective covers.
Section~\ref{sec:analytic-abundance} records relative analytic abundance
near compact subsets of the base for projective morphisms.
A reader using only the finite-data comparison can start with
Theorem~\ref{common:finite-data-frobenius} and its proof in
Section~\ref{fr:section}; neither assumes the counterexample geometry.

\subsection{Conventions}

A variety is integral. Canonical divisors are chosen compatibly on
birational models. For a divisor \(E\) on a smooth model
\(p:W\to X\), our log discrepancy is
\[
 a(E;X,B)=1+\ord_E\bigl(K_W-p^*(K_X+B)\bigr).
\]
Log canonicity means nonnegative log discrepancies, and klt means
strictly positive log discrepancies. We use the usual dlt and slc
conventions. For non-nef pseudo-effective divisors, statements
about numerical dimension use Nakayama's \(\kappa_\sigma\), as
specified in the relevant reduction. These notions are not
interchanged without a hypothesis that justifies doing so.

For nonempty systems, the Iitaka dimension is the maximum dimension of
their images. Numerical and rational-linear equivalence
are denoted by \(\equiv\) and \(\sim_{\Q}\), respectively. Unless a
different base field is specified, the constructions are over \(\C\).
A very general point avoids a countable union of proper closed
subvarieties. The transfer to arbitrary algebraically closed
characteristic-zero fields is made only after the complex argument
has been completed.
